\section{Methodology}
\label{sec:methodology}

Our goal is to quantify the correlation between static analysis metrics and functional correctness for LLM-generated code. We describe our experimental design, statistical framework, and the rationale behind key choices.

\subsection{Overview}

We compute point-biserial correlation between SA metrics (continuous) and pass@1 outcomes (binary) on HumanEval and MBPP benchmarks. To isolate the SA signal from confounding factors, we compute partial correlations controlling for code length (LOC). The core question: after accounting for the fact that shorter code tends to be both cleaner and more correct, does SA signal provide additional predictive power?

\textbf{Rationale:} Direct correlation measurement (rather than feedback loop improvement) enables rejection sampling---filtering $k$ candidates by SA score without additional LLM calls. This requires knowing whether SA predicts correctness, not whether SA feedback improves correctness.

\subsection{Static Analysis Metrics}

We extract three SA metrics capturing different quality dimensions:

\textbf{Pylint Score (0--10):} Measures code style, conventions, and potential errors. Pylint applies $\sim$400 rules encoding decades of Python best practices. Higher scores indicate cleaner code. We hypothesize style quality correlates with logical correctness because both reflect underlying code clarity.

\textbf{Mypy Error Count:} Static type checking. More type errors may indicate logical inconsistencies, though the relationship is less direct than style metrics. We include mypy to test whether type-level analysis adds predictive signal.

\textbf{Radon Cyclomatic Complexity:} Measures decision branches (if/else, loops, boolean operators). Higher complexity correlates with more potential logical errors---fewer branches mean fewer opportunities for mistakes. We expect negative correlation with correctness.

\textbf{Implementation:} We wrap each tool in subprocess calls with 30-second timeouts, following patterns validated in preliminary experiments. All tools produce valid outputs on 100\% of samples (no crashes or timeouts), as verified in our tool coverage experiment.

\subsection{Confound Control: Code Length}

Code length (LOC) is a potential confound: shorter code may be both (1) cleaner (fewer SA warnings) and (2) more correct (simpler logic). If we observe SA-correctness correlation without controlling for LOC, the relationship might be spurious---driven by length rather than SA signal.

\textbf{Partial Correlation:} We compute partial correlation controlling for LOC. If partial correlation remains strong (comparable to raw correlation), the SA signal is genuine.

\textbf{Implementation:} We use \texttt{pingouin.partial\_corr()} for partial correlation computation, with \texttt{scipy.stats.pointbiserialr()} for raw point-biserial correlation.

\subsection{Statistical Framework}

\textbf{Point-Biserial Correlation:} Appropriate when correlating a continuous variable (SA metric) with a binary outcome (pass/fail). Equivalent to Pearson correlation when one variable is dichotomous.

\textbf{Significance Testing:} We use $\alpha=0.05$ as significance threshold. Given sample sizes $N>400$, we have $>$99\% power to detect $r\geq0.20$ effects.

\textbf{Success Criterion:} We set $r\geq0.35$ as the threshold for ``moderate'' correlation, following conventional interpretation. This threshold indicates SA scores would provide meaningful signal for rejection sampling.

\subsection{Dataset}

We combine HumanEval (164 problems) and MBPP sanitized (427 problems) for 591 unique problem-solution pairs. This combination provides:

\begin{itemize}
\item \textbf{Diversity:} HumanEval emphasizes algorithmic problems; MBPP includes more practical tasks
\item \textbf{Standard benchmarks:} Both are widely used, enabling comparison with prior work
\item \textbf{Known ground truth:} Test suites define correctness unambiguously
\end{itemize}

For correlation analysis, we use canonical solutions (one completion per problem) to establish the SA-correctness relationship on known-correct code. Cross-model analysis uses LLM-generated completions to test generalization.
